\subsection{TPD 谱的三层信息：峰数、峰温与峰面积}
\label{sec:ch05b-read}

\begin{probout}
分片 a 建立了 TPD 的动力学基础与三种吸附等温式，本节不再重复，转而处理实际操作中最容易被跳过的一步——
\emph{把一张 TPD 谱读成结论}。一张谱图只提供三条可以独立读取的信息：峰数、峰温、峰面积。
本节回答：这三条信息各自指向什么结论、结论成立需要什么前提、以及在什么条件下这个指向会彻底失效。
判读的目标不是「数出几个峰就报几个物种」，而是给出「观测—前提—推断」三段式结论。
\end{probout}

\begin{principle}
\textbf{峰面积给出脱附量。} 检测器信号 $A(T)$ 与载气中脱附物种的浓度成正比。设载气总摩尔流率为 $F$
（mol\,s$^{-1}$）、标定因子 $k_A$（使浓度 $c=A/k_A$）、升温速率 $\beta=\mathrm{d}T/\mathrm{d}t$、样品质量 $m$，
则单位质量样品的脱附量为
\begin{equation}
\frac{n_{\mathrm{des}}}{m}
=\frac{F}{k_A\,\beta\,m}\int_{T_1}^{T_2} A(T)\,\mathrm{d}T .
\label{eq:ch05b-area}
\end{equation}
式中 $F/(k_A\beta m)$ 是几个量拧在一起的总标定常数：实践中不逐个测定，而是用已知量的标样整体标定
（\S\ref{sec:ch05b-quant}）。积分窗口 $[T_1,T_2]$ 必须等信号回到基线之后才关闭；窗口开得不够宽，
高温尾部被切掉，脱附量会系统性偏低，而这一点在谱图上完全看不出来。

\textbf{峰温给出脱附活化能。} 对一级脱附，Polanyi--Wigner 速率式
$\mathrm{d}\theta/\mathrm{d}t=-\nu\theta\exp[-E_{\mathrm d}/(RT)]$ 在峰顶满足
$\mathrm{d}^{2}\theta/\mathrm{d}T^{2}=0$，用 $\beta$ 换元后得
\begin{equation}
\frac{E_{\mathrm d}}{R T_{\mathrm p}^{2}}
=\frac{\nu}{\beta}\exp\!\left(-\frac{E_{\mathrm d}}{R T_{\mathrm p}}\right)
\quad\Longrightarrow\quad
\frac{E_{\mathrm d}}{R T_{\mathrm p}}
=\ln\frac{\nu T_{\mathrm p}}{\beta}-3.64 ,
\label{eq:ch05b-redhead}
\end{equation}
即通常所说的 Redhead 式。右端的 $3.64$ 来自 $\ln[E_{\mathrm d}/(RT_{\mathrm p})]$ 在常见参数区间内的近似取值，
所以 \equref{eq:ch05b-redhead} 是\emph{近似式}，而且 $\nu$ 必须事先指定（常用 $10^{13}\,\mathrm{s^{-1}}$）。
对 \equref{eq:ch05b-redhead} 求微分，得到两条后面反复用到的灵敏度关系：
\begin{equation}
\frac{\partial T_{\mathrm p}}{\partial\ln\beta}=\frac{R T_{\mathrm p}^{2}}{E_{\mathrm d}},
\qquad
\delta E_{\mathrm d}=R T_{\mathrm p}\,\delta\ln\nu .
\label{eq:ch05b-sens}
\end{equation}
\end{principle}

\begin{measure}
取 $E_{\mathrm d}=\qty{120}{kJ.mol^{-1}}$、$\nu=10^{13}\,\mathrm{s^{-1}}$，数值求解
\equref{eq:ch05b-redhead} 得 $T_{\mathrm p}\approx 423$ K（$\beta=10$ K\,min$^{-1}$）。
由 \equref{eq:ch05b-sens}：

\begin{itemize}
  \item 升温速率从 10 提高到 40 K\,min$^{-1}$（$\ln 4=1.386$），
        $\Delta T_{\mathrm p}\approx(8.314\times423^{2}/120000)\times1.386\approx 17$ K。
        也就是说，两个实验室只要升温速率不同，同一催化剂也可以差出十几开——这属于仪器差异，不是化学差异。
  \item 峰温测量误差 $\pm5$ K 只带来 $\pm1.4$ kJ\,mol$^{-1}$ 的 $E_{\mathrm d}$ 误差，可以忽略；
        真正的不确定度来自 $\nu$：$\nu$ 从 $10^{12}$ 改到 $10^{14}$（$\delta\ln\nu=4.6$）会使
        $E_{\mathrm d}$ 变动 $8.314\times423\times4.6\approx\qty{16}{kJ.mol^{-1}}$，
        即单张 TPD 谱给出的 $E_{\mathrm d}$ 精度大约只有 $\pm10\%$，而这还是在假定模型正确的前提下。
\end{itemize}

因此本节后续一律要求：\emph{报告 $E_{\mathrm d}$ 时必须同时报告所用的 $\nu$、升温速率和是否做过变速率校正}；
只报一个数而不报条件，等于没有报。
\end{measure}

要求 1 的核心结论由图~\ref{fig:ch05b-tpd-model} 概括：升温速率与再吸附都不改变脱附量，
但它们强烈改变峰形，所以\emph{峰高、峰宽、峰温都不能当作吸附量的量度，只有峰面积可以}。

\begin{figure}[htbp]
\centering
\begin{tikzpicture}
\tikzset{declare function={
  fA(\t) = \t < 423 ? exp(-((\t-423)/58)*((\t-423)/58)) : exp(-(\t-423)/95);
  fB(\t) = \t < 440 ? exp(-((\t-440)/58)*((\t-440)/58)) : exp(-(\t-440)/95);
  fC(\t) = \t < 470 ? 0.62*exp(-((\t-470)/80)*((\t-470)/80)) : 0.62*exp(-(\t-470)/175);
}}
\begin{axis}[
  width=\onecolwidth,
  height=62mm,
  xmin=280, xmax=720, ymin=0, ymax=1.22,
  xlabel={温度 $T$ / K},
  ylabel={归一化脱附速率 $r/r_{\max}$},
  axis lines=left,
  tick label style={font={\fontsize{8}{9.6}\selectfont}},
  label style={font={\fontsize{8}{9.6}\selectfont}},
  legend style={font={\fontsize{8}{9.6}\selectfont}, draw=none,
                at={(0.02,0.99)}, anchor=north west, cells={anchor=west}},
  legend cell align=left,
]
\addplot[domain=280:720, samples=320, solid, thick] {fA(x)};
\addlegendentry{$\beta=10$ K\,min$^{-1}$，无再吸附}
\addplot[domain=280:720, samples=320, dashed, thick] {fB(x)};
\addlegendentry{$\beta=40$ K\,min$^{-1}$，无再吸附}
\addplot[domain=280:720, samples=320, dotted, very thick] {fC(x)};
\addlegendentry{$\beta=10$ K\,min$^{-1}$，强再吸附}
\node[font={\fontsize{8}{9.6}\selectfont}, anchor=north west, align=left,
      text width=40mm] at (axis cs:585,1.12)
  {峰序上移 17 K：纯升温速率效应\\
   峰宽化、拖长高温尾、峰高降至 0.62：再吸附效应\\
   三条曲线面积近似相等};
\end{axis}
\end{tikzpicture}
\caption{再吸附与升温速率对一级 TPD 峰形的模型影响。读者应看三件事：虚线相对实线的峰温上移约 17 K，
与 \equref{eq:ch05b-sens} 给出的 $\Delta T_{\mathrm p}\approx(RT_{\mathrm p}^{2}/E_{\mathrm d})\ln(\beta_2/\beta_1)$
一致；点线在同样升温速率下因再吸附而峰温更高、峰宽更大、高温拖尾更长；三条曲线的峰面积近似相等，
故峰高与峰宽都不能当作吸附量的量度。局限：曲线为唯象模型（前段用高斯型上升沿，后段用指数尾），
参数是设定的，只用于说明判读规则，不代表任何具体体系。\modelmark
\figsrc{据本书 \S5.4 的判读逻辑自行构造的模型计算，非原书插图}}
\label{fig:ch05b-tpd-model}
\end{figure}

\begin{pitfall}
\textbf{峰温不是脱附活化能的直接读数。} 把 $T_{\mathrm p}$ 代入 \equref{eq:ch05b-redhead} 得到 $E_{\mathrm d}$，
这一步隐含了四个前提：（i）脱附为一级且 $\nu$ 值取定；（ii）不存在再吸附；（iii）不存在孔内扩散限制；
（iv）升温速率在峰区恒定。多孔催化剂（分子筛、氧化铝、二氧化铈）中的碱性探针分子极易在孔道内再吸附，
后果是峰温整体上移、峰形宽化、$T_{\mathrm p}$ 对 $\beta$ 的依赖被削弱。此时按 Redhead 式反算出来的
是「脱附 + 再吸附 + 扩散」的\emph{表观}活化能，可以比真实 $E_{\mathrm d}$ 高出几十 kJ\,mol$^{-1}$，
而且不同样品因孔结构不同，表观值的偏差方向还不一致，跨样品比较会得出相反的顺序。
\end{pitfall}

\begin{measure}
抑制再吸附的四条可操作手段，按优先级排列：
（i）提高载气流速，缩短脱附分子在床层中的停留时间；
（ii）减小催化剂粒径并稀释装样，缩短孔内扩散路径；
（iii）减小装样量（毫克级），缩短床层长度；
（iv）在若干流速、若干升温速率下各测一遍，把 $T_{\mathrm p}$ 对流速和 $\beta$ 外推，
检查是否收敛到同一个值——收敛才是「本征脱附」的证据。
若四条做完 $T_{\mathrm p}$ 仍在移动，就应放弃从峰温求 $E_{\mathrm d}$，只把峰温当作相对强弱的半定量指标。
\end{measure}

\begin{dataflow}
\textbf{峰数 $\to$ 位点种类。} 这是三层信息里最弱的一层，只能给出不等式：可分辨峰数 $\le$ 位点（吸附态）种类数。
把峰数直接当成物种数，需要同时满足四条：（i）每类位点只与探针分子形成一种吸附态；
（ii）不同位点的峰在温度轴上分得开（峰间距大于两倍峰宽）；（iii）吸附质在升温过程中不发生转化、
聚合或分解；（iv）不存在再吸附造成的峰间串扰。任意一条不成立，峰数与物种数就脱钩——
这一点在含硼分子筛的乙烯/丁烯 TPD 上有直接的实验证据（\S\ref{sec:ch05b-probe}）。

\textbf{峰温 $\to$ 相对强弱。} 同一系列样品、同一仪器、同一升温速率下比较峰温是可靠的半定量判据；
跨实验室、跨升温速率比较则必须先换算（\equref{eq:ch05b-sens}）。

\textbf{峰面积 $\to$ 吸附量（分峰后$\to$各类位点的量）。} 这是三层信息里最可靠的一层，
前提是基线扣除正确、积分窗口完整、标样已定标。总酸量取所有峰面积之和，
各类酸量取分峰后的面积，二者都必须给出误差范围。
\end{dataflow}

\begin{crosscheck}
TPD 给出的是「脱附出来多少探针分子」，它无法回答「脱附出来的还是不是原来的探针分子」。
原文对此有一句必须原样记住的告诫：热导检测到的信号是脱附物质热量的信号总和，
\emph{并不能说明脱出物种的化学组成}。因此凡涉及可能发生转化的体系（烯烃、醇、CO），
必须用质谱多通道同时监测（例如 $m/e=28$、$32$、$44$ 同测）或用冷阱收集后以 FID 分析产物分布，
再与 TPD 结论交叉验证。
\end{crosscheck}

\subsection{读谱实例：\texorpdfstring{O$_2$}{O2}-TPD 上的氧物种与 CO-TPD 上的吸附态}
\label{sec:ch05b-example}

以掺 Ag 的氧化锰八面体分子筛（OMS-2）为例，说明同一张谱图上「观测」与「推断」如何分开写。
该体系的制备与表征见原文 5.4.3 节。\src{上册 \S5.4.3，书内 p.325--327 / PDF p.344--346}

\begin{example}[\texorpdfstring{O$_2$}{O2}-TPD 与 CO-TPD 的观测与推断]
\begin{obs}
（1）\ce{OMS-2} 样品在 $m/e=32$ 通道上，随温度升高出现三个明显的脱附峰，
可读出的峰温为 769 K 与 861 K（第三个峰温在 OCR 层缺失），其中低温端峰 I 的峰强最小。
（2）掺入 Ag 后各脱附峰位置类似，但峰温均有不同程度降低；随 Ag 含量增加，峰 II 峰温下降最明显。
（3）CO-TPD 中同时监测 $m/e=28$、$32$、$44$：各催化剂表面几乎不出现 $m/e=28$ 信号峰，
脱附产物主要以 $m/e=44$ 出现。
（4）\ce{OMS-2} 表面只有一个较强的 \ce{CO2} 脱附信号峰；\ce{Ag-OMS-2} 出现三个 \ce{CO2} 脱附信号峰。
\end{obs}
\begin{infer}
由（1）：样品上存在三种氧物种，按原文归属，低温峰 I 为表面吸附氧，高温峰 II、III 分别为
氧化锰的表层晶格氧与体相晶格氧。\emph{核对说明}：已对照原页面图像（书内 p.325 / PDF p.344）。图 5-14 中 OMS-2（曲线 a）三个 $m/e=32$ 峰的峰温依次为 769 K、861 K、1011 K，峰 I、II、III 的标注与正文一致；0.1\% Ag/OMS-2（b）为 771、848、1000 K，0.5\% Ag/OMS-2（c）为 756、831、989 K。OCR 层缺失的第三个峰温为 1011 K，本节数值已核对一致。
由（2）：Ag 增大了表面晶格氧的扩散速率，使其更易迁移并脱附，因而比未掺 Ag 的材料具有更强的氧化性能。
由（3）：CO 在表面发生了化学吸附并与表面晶格氧反应，以 \ce{CO2} 形式脱附，
说明掺杂 Ag 的氧化锰表面 CO 氧化遵循氧化还原（Mars--van Krevelen 型）机理，而非分子态 CO 直接脱附。
由（4）：Ag 的掺入促进了新的 CO 吸附态生成；新吸附态脱附面积更大、脱附温度更低，
因而能更好地活化 C--O 键，促进 CO 氧化。
\end{infer}
\end{example}

\begin{pitfall}
把（1）读成「三种氧物种」还需要一个前置排除：$m/e=32$ 通道必须确实只来自 \ce{O2}。
若体系中同时有 \ce{CO2} 大量脱出，电子轰击离子源会由 \ce{CO2} 产生 $m/e=28$ 与 $m/e=44$ 的碎片，
但不产生 $m/e=32$，所以本例中 32 通道是干净的；换成含硫或含氯体系就必须重新逐通道核对碎片贡献。
另外要注意 TPD 本身会改变样品：原文指出随氧从表面不断迁移和脱附，材料的钾锰矿结构逐步遭到破坏、
孔道塌陷并发生烧结——所以 TPD 是\emph{破坏性}表征，同一样品不能先用它测完再去做其他表征。
\end{pitfall}

\subsection{多探针 TPD 区分酸位类型：\texorpdfstring{NH$_3$}{NH3}、乙烯与丁烯}
\label{sec:ch05b-probe}

\begin{probout}
固体酸性的常规做法是用碱性气体（\ce{NH3}、吡啶等）的 TPD 求总酸量和酸强度分布。
但单探针只能给出「强/弱」这一条轴。若改用\emph{碱强度与分子尺寸都不同}的多种探针，
同时比较它们的吸附量比例，就可以把酸位\emph{类型}也分开。本节用含硼分子筛的例子说明这条思路。
\src{上册 \S5.4.1，书内 p.320--324 / PDF p.339--343}
\end{probout}

\begin{principle}
酸性位上的吸附从强酸位开始、逐步向弱酸位发展；脱附次序相反，弱酸位上的碱性分子先脱附，
强酸位上的后脱附。因此对给定催化剂，选择合适的碱性探针并测定其脱附量，
即可得到总酸量；分峰后的峰面积给出各类酸位的酸量。当不同探针分子的碱强度存在明确顺序时，
吸附量的\textbf{比值}就成为区分位点类型的判据——这是本节真正的技术内核，
比单纯比较峰温稳健得多。
\end{principle}

\begin{dataflow}
原文用 \ce{NH3}、乙烯（\ce{C2H4}）与 1-丁烯（1-\ce{C4H8}）三种探针，在 B 与 Al 取代的 ZSM-5
和 ZSM-11 系列分子筛上做 TPD，用吸附量的相对大小反推位点类型，得到三种吸附位的模型，
并据此把样品分类。原文给出的探针碱强度顺序为 \ce{NH3} $>$ 1-\ce{C4H8} $>$ \ce{C2H4}。
\emph{核对说明}：已对照原页面图像（书内 p.320 / PDF p.339）。原书该句印作「这三种吸附质的碱性强度为 NH\textsubscript{3}>1-C\textsubscript{2}H\textsubscript{4}>C\textsubscript{4}H\textsubscript{8}」——「1-」错置在乙烯上（乙烯无位置异构），丁烯反而缺位置标号；且该顺序与同节 p.323 图 5-11 三种吸附位模型所隐含的强弱（\ce{NH3} 最强、丁烯次之、乙烯最弱）相矛盾，属原书自身问题。正文按自洽顺序写作 \ce{NH3} $>$ 1-\ce{C4H8} $>$ \ce{C2H4}。

关键观察与对应的分类结论如下：

\begin{itemize}
  \item \textbf{\ce{NH3}-TPD}：所有样品在 438 K 处均有一个脱附峰；样品 D 在 635 K 高温处另有一个较弱的脱附峰。
        样品 B\textsubscript{1}、B\textsubscript{2}、C\textsubscript{1}、C\textsubscript{2} 观察不到高温峰——
        原文归因于这些样品中 Al 含量太低。
  \item \textbf{1-\ce{C4H8}-TPD}：样品 A 与样品 Na-D（D 经钠交换）没有丁烯脱附峰；
        B\textsubscript{1}、B\textsubscript{2}、C\textsubscript{1}、C\textsubscript{2} 的丁烯峰强度随硼含量递增；
        B\textsubscript{1}、B\textsubscript{2} 只有一个峰，C\textsubscript{1}、C\textsubscript{2} 有两个峰，
        而样品 D 的两个峰明显向低温移动。
  \item \textbf{\ce{C2H4}-TPD}：样品 A 与 Na-D 没有脱附峰，B\textsubscript{1}、B\textsubscript{2} 也没有任何脱附峰，
        而 C\textsubscript{1}、C\textsubscript{2} 与 D 有两个脱附峰。
\end{itemize}

\textbf{吸附量与骨架取代数的比较}给出决定性证据：\ce{NH3} 吸附量对应于 B 和 Al 元素取代的骨架数；
而乙烯、丁烯则不同——样品 A 和 Na-D 不吸附丁烯；B\textsubscript{1}、B\textsubscript{2} 的丁烯吸附量
明显低于 \ce{NH3} 吸附量；C\textsubscript{1}、C\textsubscript{2} 的丁烯吸附量却\emph{大于} Al 和 B 的化学计量数；
样品 D 的丁烯吸附量是 \ce{NH3} 的 5 倍。原文据此判断丁烯在分子筛表面发生了聚合，
每个酸位吸附了约 5 个丁烯分子；乙烯可作类似解释。

据此可以建立三种吸附位的模型：吸附位 O 只吸附 \ce{NH3}；吸附位 I 吸附 \ce{NH3} 和丁烯但不吸附乙烯；
吸附位 II 对三者都吸附。样品 A 不吸附丁烯和乙烯，故只有吸附位 O；B\textsubscript{1}、B\textsubscript{2}
既有 O 又有 I、但没有 II；C\textsubscript{1}、C\textsubscript{2} 因丁烯吸附量远大于 \ce{NH3} 吸附量，
被认为没有吸附位 O。结论：\emph{采用碱强度不同的吸附质，可表征催化剂表面酸位的种类}。
\end{dataflow}

\begin{pitfall}
上例中丁烯吸附量「大于化学计量数」这一条本身就是误判的入口：如果按「一个酸位吸附一个分子」
去算酸量，会得到比 \ce{NH3} 滴定高数倍的酸量，进而推出「样品 C\textsubscript{1} 的酸位数远多于 C\textsubscript{2}」
这种错误结论。正确的处理是：先用 \ce{NH3} 定酸位总数，再用其他探针的\emph{比值}判断位点类型，
而不是把每种探针的吸附量都当成酸量。

更要紧的是，热导检测到的信号是脱附物质热量的信号总和，并不能说明脱出物种的化学组成。
原文把 C\textsubscript{2} 样品的脱附产物用冷阱收集后做 FID 分析，发现脱出物中除原始吸附物外
还有许多其他物质，且相对分子质量均大于乙烯，从而证实乙烯在脱附过程中于表面发生了聚合。
如果只看热导峰面积，这批数据会被完整地误读为「乙烯在强酸位上的脱附」。
\end{pitfall}

\begin{example}[峰温迁移与峰面积变化的分离]
\begin{obs}
MCM-49 经硝酸回流或水蒸气两种方式脱铝后做 \ce{NH3}-TPD.
未改性样品的谱线上可明显看到两个脱附峰：低温峰位于 \qty{219}{\degC}，高温峰位于 \qty{365}{\degC}。
硝酸脱铝后，两个峰的位置基本不变，但两峰面积均下降，且随酸浸温度提高与酸浸时间延长，下降更明显。
水蒸气脱铝后，总酸量下降，其中强酸峰面积的下降幅度大于弱酸峰；当水蒸气处理温度提高到 \qty{600}{\degC} 时，
弱酸峰与强酸峰\emph{同时向低温方向迁移}。
\emph{核对说明}：已核对并修正（原页面图像，书内 p.324 / PDF p.343）。原书全段（含节标题）一律作 \textbf{MCM-49}，「MCM-9」系 OCR 层误识，正文不再保留该异名。图 5-12 的样品名原书印为 (a) HMCM-49；(b) D-85-3；(c) D25-5；(d) D85-5；(e) D85-10，是硝酸回流脱铝这一组自身的样品系列；图 5-13（书内 p.325）另印为 (a) HMCM-49；(b) SD-550-3；(c) SD550-6；(d) SD600-3，属水蒸气脱铝的另一组，二者与原书 \S5.4.1 的样品 D（\ce{NH3}/\ce{C2H4}/1-\ce{C4H8} 探针实验中的 [Al, B]-ZSM-5）\emph{并无对应关系}——把 D-85-3、D25-5、D85 与样品 D 联系起来是 OCR 层把两节样品名混同所致。
\end{obs}
\begin{infer}
峰位不变 + 峰面积下降 $\Rightarrow$ 酸\emph{强度分布}未变，酸\emph{量}减少（硝酸脱铝）。
峰位向低温迁移 $\Rightarrow$ 酸中心\emph{强度}降低（高温水蒸气脱铝）。
两者是两条独立的结论：只报总酸量下降，会把「酸量减少」和「酸强度降低」两种完全不同的改性效果混为一谈，
而它们对催化性能的影响方向常常相反。
\end{infer}
\end{example}

原文给出的催化性能对照进一步支持这一区分：未改性 MCM-49 在实验条件下丙烯转化率 99.5\%、
异丙苯选择性 73.6\%；脱铝改性降低了烷基化反应活性和异丙苯选择性，但提高了收率；
与脱铝前相比，常温下经硝酸处理 5 h 的脱铝 MCM-49 在保持催化活性相当的情况下，异丙苯收率提高了 5.3\%。

\subsection{TPR 谱的解析：峰温、耗氢量与物种归属}
\label{sec:ch05b-tpr}

\begin{probout}
TPR 的原始输出是「耗氢速率随温度的变化」。要把它变成物种信息，需要三样东西同时到位：
峰温（给出还原的\emph{相对}难易）、峰面积（给出耗氢\emph{量}，进而给出平均价态变化）、
以及一个能够把峰与物种对应起来的外部判据。本节用氧化物催化剂的两个实例说明怎么做到，
并给出「峰温不能单独作为判据」的具体理由。
\src{上册 \S5.5.1--\S5.5.4，书内 p.327--335 / PDF p.346--354}
\end{probout}

\begin{dataflow}
\textbf{实例一：CuO--CeO\textsubscript{2} 中三种 Cu 物种的区分。}
样品为柠檬酸溶胶--凝胶法制备的 10 mol\% CuO--CeO\textsubscript{2}：空气中 \qty{800}{\degC} 焙烧 4 h
得样品 AS；先在 \ce{N2} 中 \qty{800}{\degC} 焙烧、再在空气中 \qty{400}{\degC} 焙烧 4 h 得样品 NSA4。
部分 NSA4 与 AS 经硝酸处理以除去表面高分散 CuO 与晶相 CuO，分别得 NSA4-HI 与 AS-HI。

CO-TPR 的观测：NSA4 在 \qty{120}{\degC} 和 \qty{175}{\degC} 附近分别出现两个低温还原峰 $\alpha$ 与 $\beta$；
经硝酸处理后 $\alpha$ 峰消失，仅剩 $\beta$ 峰。AS 可观测到位于 \qty{175}{\degC} 附近的 $\beta$ 峰
以及处于 $\beta$ 峰高温方向的肩峰 $\gamma$；经硝酸处理后 $\gamma$ 峰消失，$\beta$ 峰变得更对称。

归属逻辑：\emph{硝酸处理}只洗掉表面高分散物种与晶相物种，因此「能被洗掉的峰」对应表面高分散或晶相物种。
结合 XRD 结果，原文判定 $\alpha$ 峰由表面高分散 CuO 的还原引起，$\gamma$ 峰由晶相 CuO 的还原引起，
而 $\beta$ 峰归属为进入 CeO\textsubscript{2} 晶格的 Cu$^{2+}$ 的还原。
这条链的关键在于：\textbf{峰温只用来排序，归属靠「选择性去除 + 独立结构表征」}。
单位反应速率的计算进一步给出定量结论：以 Cu 计，高分散 CuO 物种对 CO 氧化的贡献最大，
进入 CeO\textsubscript{2} 晶格的 Cu$^{2+}$ 贡献最小。
\emph{核对说明}：已对照原页面图像（书内 p.329 / PDF p.348）。表 5-6 的数值已逐项核对一致——以催化剂质量计的反应速率为 N8A4 3.83、A8 2.28、N8A4H1 与 A8H1 均为 0.40 mmol$_{\mathrm{CO}}$/(g$_{\mathrm{cat}}$\textperiodcentered h)；以 Cu 计的单位反应速率为 N8A4 101.8、A8 60.6、N8A4H1 与 A8H1 均为 21.3 mmol$_{\mathrm{CO}}$/(g$_{\mathrm{Cu}}$\textperiodcentered h)；按物种拆分的单位反应速率为 $a$（表面高分散 CuO）183.3、$c$（体相 CuO）100.4、$b$（进入 CeO\textsubscript{2} 晶格的 Cu$^{2+}$）21.3 mmol$_{\mathrm{CO}}$/(g$_{\mathrm{Cu}}$\textperiodcentered h)。原书表注把 $c$ 印作「体相 CuO」，而同页正文把对应的 $\gamma$ 峰说成「晶相 CuO」的还原，同一物种两种叫法，属原书自身问题。

\textbf{实例二：Ce\textsubscript{x}Ti\textsubscript{1-x}O\textsubscript{2} 固溶体的双峰结构。}
溶胶--凝胶法制备的系列固溶体做 H\textsubscript{2}-TPR：对 $x>0.3$ 的样品，
分别在 \qty{650}{\degC} 与 \qty{850}{\degC} 处出现两个还原峰；对富 Ti 端的
Ce\textsubscript{0.1}Ti\textsubscript{0.9}O\textsubscript{2} 样品只有一个互相重叠的还原峰；
TiO\textsubscript{2} 在该温度范围内无还原。

可读出的规律有四条：（i）\qty{650}{\degC} 的低温还原峰随 TiO\textsubscript{2} 含量增加逐渐向高温方向移动，
而高温还原峰无明显变化；（ii）低温还原峰强度随 TiO\textsubscript{2} 含量增加而加强，并在 $x=0.4$ 时达到最大值；
（iii）高温还原峰的变化较复杂，$x$ 从 1 降至 0.6 时该峰强度下降，在 $x=0.5$ 时达到最大值；
（iv）从系列固溶体还原峰的相对面积看，\qty{650}{\degC} 处峰面积的增加主要来自 TiO\textsubscript{2} 含量的增加，
而 \qty{850}{\degC} 处出现的尖锐还原峰可能与样品中出现的立方相有关。

耗氢量数据给出的结论是：Ce\textsubscript{x}Ti\textsubscript{1-x}O\textsubscript{2} 复合氧化物的耗氢量
远高于纯 CeO\textsubscript{2}，表明 TiO\textsubscript{2} 的加入促进了 CeO\textsubscript{2} 的还原；
$x=0.2\sim0.6$ 时还原样品中 Ce 的平均价态明显低于纯 CeO\textsubscript{2}，
即复合氧化物具有比 CeO\textsubscript{2} 更高的储氧能力。
\emph{核对说明}：已对照原页面图像（书内 p.330--331 / PDF p.349--350）。表 5-7 的耗氢量（mmol/g）与括号内误差已核对一致：$x=0$ 为 0；0.1 为 0.356（0.016）；0.2 为 1.072（0.050）；0.3 为 1.250（0.062）；0.4 为 1.920（0.096）；0.5 为 2.144（0.108）；0.6 为 1.964（0.098）；0.7 为 1.830（0.092）；0.8 为 1.428（0.072）；0.9 为 1.384（0.070）；1.0 为 0.580（0.030）。表注定义 $b$ 为混合氧化物总的还原程度、$c$ 为混合氧化物 CeO\textsubscript{2} 的还原程度。原书正文印作「复合氧化物中 Ce 的平均价态为 ＋3，远高于纯 CeO\textsubscript{2} 中 Ce 的平均价态 ＋3.8」，「远高于」与其自身印出的两个数值（＋3 低于 ＋3.8）相矛盾，属原书自身问题；本处按数值自洽取「明显低于」。
\end{dataflow}

\begin{measure}
\textbf{再氧化实验是判断「可逆性」的标准做法。}
把已还原样品在不同温度（\qty{500}{\degC}、\qty{300}{\degC}、\qty{100}{\degC}）再氧化后重做 TPR：
与新鲜样品相比还原峰变弱。对 Ce\textsubscript{0.8}Ti\textsubscript{0.2}O\textsubscript{2}，
第二次还原的峰位置没有改变；而对 Ce\textsubscript{0.4}Ti\textsubscript{0.6}O\textsubscript{2}，
第二次还原峰位置向低温方向偏移。按峰面积比较，前者在 \qty{500}{\degC}、\qty{300}{\degC}、\qty{100}{\degC}
再氧化后的耗氢量分别为新鲜样品的 98\%、80\% 和 80\%，后者为 90\% 和 78\%，说明前者更容易被再次氧化。
\emph{核对说明}：已对照原页面图像（书内 p.331 / PDF p.350）。原文给出的比例是：Ce\textsubscript{0.8}Ti\textsubscript{0.2}O\textsubscript{2} 在 \qty{500}{\degC}、\qty{300}{\degC}、\qty{100}{\degC} 再氧化后，耗氢量为新鲜样品的 98\%、80\% 和 80\%，三个温度点齐全；Ce\textsubscript{0.4}Ti\textsubscript{0.6}O\textsubscript{2} 只给出 90\% 和 78\% 两个比例，其余温度点的比例原书未给出。原书亦未给出这两组比例的误差范围。

这类实验的判读要点是：\emph{峰位变化}说明还原物种的性质变了（例如体相氧的可动性改变），
\emph{峰面积变化}说明可还原的量变了。两者要分开报告，不能合并成一句「还原性变弱」。
\end{measure}

\begin{pitfall}
\textbf{峰数 $\ne$ 物种数，在氧化物 TPR 上尤其危险。}
Ce\textsubscript{0.1}Ti\textsubscript{0.9}O\textsubscript{2} 只有一个「互相重叠」的还原峰，
如果按峰数判物种数就会得出「只有一个物种」的结论，而实际上该峰内至少包含两部分的贡献；
反过来，CuO--CeO\textsubscript{2} 中 AS 样品的 $\gamma$ 只是 $\beta$ 峰高温方向的\emph{肩峰}，
把它当作独立物种并分别积分，会同时得到错误的物种数和错误的定量结果。
可靠的峰数判定至少需要三条中的两条同时成立：分峰拟合与二阶导数分析、
变升温速率下各分量是否按同一规律移动、以及独立变量实验（酸洗、负载量序列、XRD 物相）是否支持该归属。
\end{pitfall}

\begin{crosscheck}
CO-TPR 与 H\textsubscript{2}-TPR 的结果可以互为校验，但两者不可互换。PdO/CeO\textsubscript{2} 体系说明了原因：
催化剂表面可能存在金属 Pd，而金属 Pd 有较强的氢溢流作用，\ce{H2} 在金属上解离后溢流到载体，
使载体氧化物在低温被还原——因此单靠 H\textsubscript{2}-TPR 无法区别催化剂中的晶相 PdO 与高分散 PdO。
CO 不能解离，溢流比氢困难得多，所以改用 CO 作还原剂做 CO-TPR 能得到更丰富的还原物种信息。
原文的观测是：实验温度范围内未发现 CeO\textsubscript{2} 消耗 CO 的信号，而在 \qty{320}{\degC} 处
出现一个 CO 消耗峰；用 CO 作还原剂时 Pd 催化剂有三个 CO 消耗峰（记为 $\alpha$、$\beta$、$\gamma$），
而用 \ce{H2} 作还原剂时只有一个还原峰。
\emph{核对说明}：已核对并修正（原页面图像，书内 p.333 / PDF p.352）。原书此处全部作 \ce{CH4}：「由于 \ce{CH4} 氧化活性与 $\alpha$ 峰不存在任何对应关系，故 $\alpha$ 物种不可能是 \ce{CH4} 氧化活性中心；而 $\beta$ 物种和 $\gamma$ 物种可能是 \ce{CH4} 氧化活性中心，因为 $\beta$、$\gamma$ 峰变化与 \ce{CH4} 氧化活性有一定的对应关系。所以，可以认为高分散的 PdO 是 CO 氧化的活性物种，晶粒较大的 $\beta$（或 $\gamma$）物种是 \ce{CH4} 氧化活性物种」。即原书对两个反应的归属本来就分开陈述：$\alpha$（高分散 PdO）对应 CO 氧化、$\beta$/$\gamma$ 对应 \ce{CH4} 氧化，\emph{本身并不矛盾}；此前所谓「前后不一致」源于 OCR 层把两处 \ce{CH4} 误读为 CO，正文已按原书改回。

\textbf{CO-TPR 的边界}：高温时 CO 容易发生歧化反应而生成 \ce{CO2}，这一点在实际操作中应特别注意；
歧化产生的 \ce{CO2} 会被计入消耗信号，使耗氢（耗 CO）量偏高。
\end{crosscheck}

\subsection{为什么 TPR 峰温不能单独作为还原难易的判据}
\label{sec:ch05b-tpr-limit}

\begin{probout}
「峰温低 = 容易还原」是一个过于粗糙的经验规则。本节说明：峰温是动力学量，
它同时受升温速率、载气流速与氢分压、样品量与粒度、氢溢流、以及还原程度本身影响；
并且在同一样品的还原与氧化两个方向上，难易顺序可以\emph{相反}。
\src{上册 \S5.5.4，书内 p.334--335 / PDF p.353--354}
\end{probout}

\begin{dataflow}
V\textsubscript{2}O\textsubscript{5}/TiO\textsubscript{2} 体系给出了最清楚的例子。
纯 V\textsubscript{2}O\textsubscript{5} 的 H\textsubscript{2}-TPR 在 910 K、932 K、971 K、1044 K 与 1102 K
（后两个未在图中显示）处出现多个还原峰，对应由 V\textsubscript{2}O\textsubscript{5} 逐步还原为
V\textsubscript{2}O\textsubscript{3} 的中间过程；结合 XRD 可以判定第一个还原峰为
V\textsubscript{2}O\textsubscript{5} 到 V\textsubscript{6}O\textsubscript{13} 的转化，第二个还原峰为
V\textsubscript{6}O\textsubscript{13} 到 \ce{VO2} 的转化，最后一个还原峰为 V\textsubscript{2}O\textsubscript{3} 的形成。\emph{核对说明}：已对照原页面图像（书内 p.334 / PDF p.353）。原书印为「对于 \ce{V2O5} 样品，在 910 K、932 K、971 K、1044 K 和 1102 K（后两个未显示），对应于由 \ce{V2O5} 转变为 \ce{V2O3} 的中间过程，通过结合 XRD 数据，可以判定第一个还原峰为 \ce{V2O5} 到 \ce{V6O13} 的转化，第二个还原峰为 \ce{V6O13} 到 \ce{VO2} 的转化，最后一个还原峰为 \ce{V2O3} 的形成」。还原峰温度序列与第二个峰的中间相对应关系（\ce{V6O13} $\rightarrow$ \ce{VO2}）均完整，OCR 层并无缺失；正文此前只引到 971 K 且略去第二个峰的归属，现已按原书补齐。

负载后行为改变：高负载催化剂 EL10V8（8 wt\% V\textsubscript{2}O\textsubscript{5}/TiO\textsubscript{2}）
有两个还原峰，比纯 V\textsubscript{2}O\textsubscript{5} 的还原温度降低了约 100 K，
说明载体 TiO\textsubscript{2} 促进了还原性能。其中 802 K 处的还原峰被认为是单层分散
V\textsubscript{2}O\textsubscript{5} 的还原，879 K 的高温还原峰则是聚合态 V\textsubscript{2}O\textsubscript{5} 的还原；
低负载催化剂 EL1V（1 wt\% V\textsubscript{2}O\textsubscript{5}/TiO\textsubscript{2}）只有一个位于 760 K 的还原峰，
归属于单层分散 V 物种的还原。

把同一批样品再做 TPO（\S\ref{sec:ch05b-tpo}）：纯 V\textsubscript{2}O\textsubscript{5} 有两个分别位于 751 K
和 843 K 的氧化峰，\emph{比 TPR 中的还原峰温度更低}；从峰型看，TPO 中的峰与 TPR 中的还原峰一一对应，
表明 V 物种的氧化与还原过程可逆。EL10V8 有一个肩峰并在 647 K 处有一个强氧化峰；
EL1V 在 668 K 处有氧化峰。

\textbf{结论（必须成对陈述）}：与纯 V\textsubscript{2}O\textsubscript{5} 相比，
负载型 V\textsubscript{2}O\textsubscript{5}/TiO\textsubscript{2} 催化剂更容易被还原\emph{也}更容易被氧化，
说明载体 TiO\textsubscript{2} 促进了氧化还原性能；进一步看，
高负载 EL10V8 中聚合态 V\textsubscript{2}O\textsubscript{5} 物种\emph{很容易被氧化但很难被还原}，
而低负载催化剂上单层分散的 V\textsubscript{2}O\textsubscript{5} 物种\emph{很容易被还原但很难被氧化}。
同一份样品在还原与氧化两个方向上给出相反的顺序——这说明「峰温低 = 这种物种更活泼」这一说法，
在不指明方向时是无意义的。
\end{dataflow}

\begin{principle}
把 TPR 峰温与物种性质联系起来的常见推理有四类，各自的成立条件不同：

\begin{itemize}
  \item \textbf{金属--载体相互作用}：相互作用越强、越容易生成难还原的界面化合物（如表面尖晶石），
        还原峰温越高。成立条件是样品经历了相同的预处理（见 \S\ref{sec:ch05b-quant}）。
  \item \textbf{双金属/合金化}：第二金属引入后峰温位移反映合金化与表面富集，
        但同样会被氢溢流和氢在金属上的强吸附所掩盖。
  \item \textbf{分散度}：高分散物种表面能高、界面氧易移除，峰温通常较低——
        \emph{但这不是普遍规律}。当金属--载体相互作用强到出现 SMSI 时，高分散反而更难还原。
        因此分散度必须由化学吸附或 XRD 线宽独立测定，不能用 TPR 峰温反推。
  \item \textbf{催化活性}：确有一系列工作显示「低温可还原性」与反应活性正相关
        （见本节末的初始耗氢速率法），但相关不等于判据，仍需 TOF 或转化率数据支持。
\end{itemize}
\end{principle}

\begin{table}[htbp]
\centering
\caption{同一样品 TPR 峰温发生位移的可能原因与判别实验。读者应先做判别实验，再决定是否把峰温位移
解释为化学差异；本表为方法学归纳。
\src{据上册 \S5.5.1--\S5.5.4，书内 p.327--335 / PDF p.346--354；\S5.5.6，书内 p.336--338 / PDF p.355--357 归纳。}}
\label{tab:ch05b-tpr-shift}
\small
\begin{tabular}{@{}p{34mm}p{50mm}p{54mm}@{}}
\toprule
观察到的位移 & 可能的物理原因 & 判别实验 \\
\midrule
提高升温速率后峰温上移 & 动力学滞后，符合
$\partial T_{\mathrm p}/\partial\ln\beta=RT_{\mathrm p}^{2}/E_{\mathrm d}$ &
变 $\beta$ 测一组，按该式外推至 $\beta\to0$ 再比较 \\
提高载气流速后峰温下移 & 再吸附或产物水被抑制 & 变流速测一组，看 $T_{\mathrm p}$ 是否收敛 \\
降低装样量后峰温下移 & 床层传热/传质滞后 & 变装样量，同时检查峰面积是否守恒 \\
加装冷阱后峰温下移、峰形变尖 & 水的陷获被消除 & 冷阱前后对比，检查基线是否更平 \\
低负载样品峰温异常低 & 高分散或氢溢流 & 用 CO-TPR 重复；查 XRF/化学吸附测分散度 \\
峰温与活性顺序不一致 & 峰温控制的不是速控步骤 & 补做 TOF 与表观活化能 \\
\bottomrule
\end{tabular}
\end{table}

\emph{核对说明}：已对照原页面图像（书内 p.313、p.332—334、p.337—338 / PDF
p.332、p.351—353、p.356—357，上册 §5.3.2、§5.5.3、§5.5.4、§5.5.6）。
逐行核对：第 5 行「低负载样品峰温异常低」有一个逐字实例——p.334 记低负载的 EL10V1（$1\%$
V$_2$O$_5$/TiO$_2$）只有一个 760 K 还原峰，归为单层分散的 V$_2$O$_5$ 物种，
而高负载 EL10V8（$8\%$）在 802 K 与 879 K 有两个峰，方向与该行相符；其判别实验「用
CO-TPR 重复」有据于 p.333：因 Pd 的氢溢流使 H$_2$-TPR 无法区分高分散与晶相 PdO，
故改用不能解离、溢流比 H$_2$ 困难的 CO 作还原剂。但 p.332 给出相反方向的实例：
PdO/CeO$_2$ 系列随负载量增加还原峰温向低温方向移动，低负载量时还原温度升高，原因是分散态 PdO 与
CeO$_2$ 相互接触并被隔离、相互作用增强而较难还原，故本行只宜作为「需逐样品判断」的提示。第 6
行「峰温与活性顺序不一致」有正反两侧实例：p.332 记 PdO/Al$_2$O$_3$ 的还原温度低于
PdO/CeO$_2$，而后者的 CO 氧化活性明显更高，原书据此明言上述还原温度与 CO
氧化活性的对应关系不能成立；p.338 记 LSMO 与 xAu/LSMO
系列的还原温度递减规律与该系列反应活性一致（还原温度越低活性越好），属本行未涵盖的相反情形。
改编补充：第 1—4 行（升温速率、载气流速、装样量、冷阱引起的峰温位移方向及相应判别实验）
在本次所读页未见给出方向。原书 p.313 只把载气流速、装量、升温速率列为影响 TPD 峰形的因素：载气流速一般
$30\sim50$ mL/min，装填量 $50\sim200$ mg、粒度 $40\sim80$ 目，
升温速率一般 $10\sim20$ $^{\circ}$C/min（过大则峰易重叠、过小则信号减弱），
并未给出峰温位移方向。书内 p.327—331（§5.5.1—§5.5.2）与 p.335—336 本次未逐页核对。


\subsection{TPR 与 TPO 的定量前提}
\label{sec:ch05b-quant}

\begin{probout}
TPR/TPO 的定量结论（耗氢量、平均价态、储氧能力）都建立在四个前提上：
标样标定、载气流速与升温速率的匹配、水的陷获、再吸附抑制。
四项缺一项，耗氢量的误差就可以超过物种归属本身带来的差异。本节给出每项的物理原因与操作对策。
\src{上册 \S5.5.6，书内 p.336--338 / PDF p.355--357}
\end{probout}

\begin{measure}
\textbf{（1）标样标定。} \equref{eq:ch05b-area} 中的总标定常数 $F/(k_A\beta m)$ 必须实验测定。
两种可行做法：其一，用定量管（六通阀）向同一套气路脉冲注入已知摩尔数的还原性气体或已知量的 \ce{H2}，
由峰面积得到 $k_A$；其二，用已知耗氢量的参比物（如已知量的 CuO）在同条件下做 TPR，
由「已知耗氢量 / 测得峰面积」确定标定系数。标定必须在\emph{与样品相同的载气组成、流速、升温程序}下进行——
TCD 的响应因子与气体组成有关，用 \ce{H2}/Ar 标定却在纯 Ar 下标定，误差可以到 10\% 以上。
原文的表 5-7 在耗氢量后附了括号内的误差，说明作者是把积分本身的不确定度一并报告的；
如果只报一个没有误差的数，读者无法判断两个样品之间的差别是否真实。

\textbf{（2）载气流速与升温速率的匹配。} 二者共同决定脱附/还原产物在床层中的停留时间与热滞后。
流速过低：产物在床层内停留时间长，再吸附与水的陷获加剧，峰温上移、峰形拖尾；
流速过高：样品床层内温度梯度增大，且信号被稀释、信噪比下降；升温速率过高：峰温按
\equref{eq:ch05b-sens} 上移、相邻峰重叠加重、样品与热电偶之间出现热滞后。
常规折中是 5\%--10\% 的 \ce{H2}/Ar、流速 $10\sim20$ mL\,min$^{-1}$、
$\beta=5\sim10$ K\,min$^{-1}$、装样 $20\sim50$ mg，稀释后装填；但\emph{一套参数不能跨体系照搬}：
分子筛等高比表面载体需要更高的流速，氧化物粉末可以低一些。判断匹配是否合适、看三点：
峰形是否对称、基线是否平、同一份样品把流速或 $\beta$ 变一倍后峰面积是否不变。

\textbf{（3）水的陷获。} H\textsubscript{2}-TPR 必然生成水，水会强烈吸附在 \ce{Al2O3}、分子筛、
CeO\textsubscript{2} 等载体上。后果有三：其一，水在 TCD 上造成的响应会被误算成耗氢；
其二，水在床层内积累会改变局部 \ce{H2} 分压，使后续还原峰位移；
其三，水脱附本身形成一个宽峰，与还原峰重叠。
对策是在样品与检测器之间加冷阱（干冰--丙酮或液氮--乙醇），并核对加冷阱前后基线；
原文在讨论 CO-TPR 时还提示高温 CO 歧化生成 \ce{CO2}，这类副产物同样需要单独考虑。

\textbf{（4）再吸附抑制。} TPD 侧的再吸附已在 \S\ref{sec:ch05b-read} 讨论；
TPR 侧的再吸附主要是水的再吸附，与（3）同源。此外，金属表面上的氢溢流可以看作一种「逆向再吸附」：
\ce{H2} 在金属上解离后迁移到载体，使本不该在该温度还原的物种被还原，
这在 Pd、Pt、Ni 催化剂上几乎无法完全避免，只能通过改用 CO 作还原剂来绕开
（\S\ref{sec:ch05b-tpr} 的 crosscheck）。
\end{measure}

\begin{pitfall}
\textbf{氧化态未预处理的 TPR 无意义。} 耗氢量对应的是「从某一确定的起始价态还原到某一确定的终态」，
如果起始价态不确定，峰面积就无法换算成化学计量。样品在空气中久置、或上一次实验后未彻底再氧化，
表面会带有不确定量的还原态物种。标准流程是：先在 \ce{O2} 或空气中于规定温度预处理，
再在惰性气氛下降温至室温，然后切入 \ce{H2}/Ar 开始程序升温。
同样的纪律适用于 TPD：\ce{NH3}-TPD 前必须先在惰性气流中于 $500\sim550\,\degC$ 脱除预吸附水与杂质，
再降温吸附 \ce{NH3}；否则谱图上第一个峰很可能只是水的脱附。

\textbf{耗氢量不扣基线。} TCD 的基线本身会随温度漂移（热导率随温度变化、水的陷获、载气中
\ce{H2} 的消耗造成的分压变化）。若不做基线扣除，漂移会被一并积分进去。
正确做法是：用同质量的纯载体（或惰性参比）在完全相同条件下走一遍空白程序得到基线，
或至少用峰前后基线的线性外推作为扣除线，并在报告中写明积分窗口的起止温度。
基线扣除把峰面积改变 10\%--20\% 是常见的，而这个量级恰好与「Cu 是高分散还是晶相」的差别同量级。
\end{pitfall}

\begin{dataflow}
\textbf{区分表面还原与体相还原的两种实用方法。}
TPR 谱图给出的往往是表面与体相还原的\emph{混合}信息，而催化反应只发生在表面，
因此「如何区分表面和体相氧化物物种」是关键问题，在普通 TPR 谱图上又比较难以实现。
原文介绍了两条经验路线：

\textbf{（i）比表面积归一化作图法。} 对一系列比表面积不同、组成相同的样品分别测 TPR，
把某个还原峰的耗氢量按「单位质量」与「单位表面积」两种方式归一化后互相作图：
$N_{\mathrm{total}}/M_{\mathrm{sample}}=a+b\,(A_{\mathrm{sample}}/M_{\mathrm{sample}})$，
其中 $N_{\mathrm{total}}$ 为某一还原峰的耗氢量、$M_{\mathrm{sample}}$ 为样品质量、
$A_{\mathrm{sample}}$ 为样品比表面积，$a$、$b$ 为常数。
若所有数据点落在一条直线上，说明该还原峰对应的是\emph{表面}氧化物的还原；
若偏离直线，则含有体相贡献。原文用它检验了一系列 CoCr\textsubscript{2}O\textsubscript{4} 尖晶石样品
（分别在 \qty{400}{\degC}、\qty{500}{\degC}、\qty{600}{\degC}、\qty{700}{\degC} 焙烧，
比表面积不同，装样量皆为 25 mg）：低温还原峰（$\alpha$ 峰）归属为催化剂中高价态 Cr 物种的还原，
但「是否属于表面 Cr 物种的还原」难以直接归属；按上述方式作图并对数据拟合后得到一条直线，
说明该还原峰确实由表面 Cr 的还原引起。
\emph{核对说明}：已对照原页面图像（书内 p.337 / PDF p.356）。经验公式原书完整印为 $\frac{N_{\mathrm{total}}}{M_{\mathrm{sample}}}=a+b\times\frac{A_{\mathrm{sample}}}{M_{\mathrm{sample}}}$，其中 $N_{\mathrm{total}}$ 为某一还原峰的耗氢量、$M_{\mathrm{sample}}$ 为样品的质量、$A_{\mathrm{sample}}$ 为样品的比表面积，$a$ 和 $b$ 为常数；以 $N_{\mathrm{total}}/M_{\mathrm{sample}}$ 对 $A_{\mathrm{sample}}/M_{\mathrm{sample}}$ 作图，若得到一条直线即说明该还原峰是表面氧化物的还原。原书对 CoCr\textsubscript{2}O\textsubscript{4}-4/-5/-6/-7（比表面积分别为 91.3、51.9、23.2、6.1 m\textsuperscript{2}/g，装样量皆为 25 mg）的低温 $\alpha$ 峰（耗氢量分别为 0.734、0.403、0.186、0.051 mmol/g）作此图并拟合，得 $Y=5.84+7.79X$（$R^2=0.99999$），是一条直线，说明该峰由表面 Cr\textsuperscript{6+} 的还原引起。公式与常数在原书中完整可辨，此前「公式与常数完全失真」的判断系 OCR 层所致。

\textbf{（ii）初始耗氢速率法。} 原文引述 A.\,T.\,Bell 等的做法：催化剂的低温还原性可用
「样品第一个还原峰中前 25\% 的氧的脱除速率」来评估。具体做法是把 TPR 曲线换算成
「初始耗氢量—温度」关系，然后在\emph{相同初始耗氢速率}（同一纵坐标值）下比较各样品所需的温度，
还原温度越低则低温可还原性越好。用这一处理比较规则三维大孔结构的 La--Sr--Mn--O 钙钛矿
（LSMO）与负载 Au 的 Au/LSMO 系列（$x$Au/LSMO，$x=3.4$--$7.9$ wt\%），
作者报告的还原温度递减规律为 6.4Au/LSMO $<$ 7.9Au/LSMO $<$ 3.4Au/LSMO $<$
LSMO $\approx$ 6.2Au/bulk LSMO $<$ bulk LSMO，且该顺序与催化反应活性顺序一致，
意味着催化剂的低温可还原性在催化反应中起到了关键作用。
\emph{核对说明}：已对照原页面图像（书内 p.338 / PDF p.357）。原书给出负载量为 $x$Au/LSMO、$x=3.4$--$7.9$ wt\%，图 5-28 图例中的样品为 LSMO、3.4Au/LSMO、6.4Au/LSMO、7.9Au/LSMO、bulk-LSMO 与 6.2Au/bulk-LSMO，各数值自洽，并不冲突。在相同初始耗氢速率（同一纵坐标值）下，原书给出的还原温度递减规律为 6.4Au/LSMO $<$ 7.9Au/LSMO $<$ 3.4Au/LSMO $<$ LSMO $\approx$ 6.2Au/bulk LSMO $<$ bulk LSMO，并指出该规律与催化反应活性（还原温度越低活性越好）一致，正文已据以补全。

\textbf{为什么必须做这一步}：直接用峰温排序会被峰高（即耗氢量）差异污染——
一个峰又高又陡的样品，其峰顶温度会天然低于一个又低又宽的样品，即使两者本征还原性相同。
初始耗氢速率法把「速率」固定在同一个值上再比温度，去掉了这一污染。
\end{dataflow}

\subsection{TPO 与积炭：三种信号的伴生关系}
\label{sec:ch05b-tpo}

\begin{probout}
TPO 用于表征催化剂的氧化性能与表面积炭。与 TPR 不同，TPO 往往要同时监测 \ce{CO2}、\ce{H2O}
与 \ce{O2} 三种信号，而\emph{判读的钥匙不在单个峰，而在峰与峰是否伴生}。
\src{上册 \S5.5.5，书内 p.335--336 / PDF p.354--355}
\end{probout}

\begin{dataflow}
样品为 13\%（质量分数）Co/\ce{Al2O3} 与 8\% Co/\ce{Al2O3}，经 \qty{400}{\degC} \ce{H2} 还原
\qty{2}{\hour} 后，通 \ce{CH4}/\ce{CO2} 混合气（原文给出 \ce{CH4}/\ce{CO2} 比例 0.93）
在 \qty{700}{\degC} 反应 \qty{1}{\hour} 积炭，随后做 TPO。

观测：13\% Co/\ce{Al2O3} 催化剂在 \qty{130}{\degC}、\qty{270}{\degC}、\qty{536}{\degC} 附近
分别出现三个 \ce{CO2} 峰；在 \qty{148}{\degC}、\qty{218}{\degC}、\qty{275}{\degC} 附近出现三个 \ce{H2O} 峰；
并在 \qty{266}{\degC} 与 \qty{540}{\degC} 附近出现两个 \ce{O2} 消耗峰。而 8\% Co/\ce{Al2O3} 催化剂的
\ce{CO2} 峰相对简单，高温峰消失。

归属逻辑：结合 \ce{CH4}/\ce{CO2} 的 TPD 实验，低温（$<\qty{150}{\degC}$）出现的 \ce{CO2} 峰
是催化剂表面预吸附 \ce{CO2} 引起的；较高温度出现的 \ce{CO2} 则来自催化剂表面积炭的氧化，
因此相应的温度上会出现耗氧峰。\ce{H2O} 峰有两个来源：
一方面反应过程中发生水煤气逆变换反应，生成的 \ce{H2O} 吸附在催化剂上、升温时脱附——
这类 \ce{H2O} 峰\textbf{不与} \ce{CO2} 脱附峰伴生，而\textbf{与} \ce{O2} 消耗峰伴生；
另一方面反应过程中催化剂表面结焦形成 \ce{C_xH_y}（$x>y$），其氧化过程会同时产生 \ce{CO2} 和 \ce{H2O}，
于是 TPO 过程中 \ce{CO2} 峰、\ce{H2O} 峰与耗 \ce{O2} 峰同时出现。

由此推出结构信息：13\% Co/\ce{Al2O3} 上有两种积炭中心，而 8\% Co/\ce{Al2O3} 上只有一种积炭
（结焦）中心；原文认为这些中心与催化剂中 Co 和 \ce{Al2O3} 不同的作用程度有关。
\emph{核对说明}：已对照原页面图像（书内 p.335 / PDF p.354）。原书该段完整印为「13\% Co/\ce{Al2O3} 催化剂在 \qty{130}{\degC}、\qty{270}{\degC} 和 \qty{536}{\degC} 附近分别出现三个 \ce{CO2} 峰；\qty{148}{\degC}、\qty{218}{\degC} 和 \qty{275}{\degC} 附近处出现三个 \ce{H2O} 峰；在 \qty{266}{\degC} 和 \qty{540}{\degC} 附近出现两个 \ce{O2} 的消耗峰」——三个 \ce{CO2} 峰温与两个 \ce{O2} 消耗峰温均在原书中给出，不存在缺失。此前正文只写「约 \qty{540}{\degC} 附近出现两个 \ce{O2} 消耗峰」，是 OCR 层失真造成的漏读，现已按原书补全首个 \ce{CO2} 峰温并补齐 \ce{O2} 峰温（另见本页正文对 \ce{CH4}/\ce{CO2} 进气的更正）。
\end{dataflow}

\begin{pitfall}
如果只测 \ce{CO2} 一条通道，上述两类 \ce{H2O} 无法区分，低温 \ce{CO2} 峰也会被误判为积炭。
更常见的误判是：把「耗氧峰」当成「积炭量」的唯一量度——耗氧量取决于积炭的 H/C 比，
而 H/C 比恰恰因结焦条件而异，因此积炭量必须由 \ce{CO2} 峰面积（碳）与 \ce{H2O} 峰面积（氢）分别定，
再用耗氧量做交叉核对。
\end{pitfall}

\subsection{TPSR：从脱附谱读出反应机理}
\label{sec:ch05b-tpsr}

\begin{probout}
程序升温表面反应（TPSR）是指在程序升温过程中同时发生表面反应和脱附。
本节回答：TPSR 与 TPD 的界限在哪里、它能给出 TPD 给不出的什么信息、
以及如何从多个产物的峰形对应关系反推速控步骤。
\src{上册 \S5.3.6 与 \S5.6，书内 p.320、p.338--340 / PDF p.339、p.357--359}
\end{probout}

\begin{principle}
TPSR 大致有两种做法。第一种：先将已预处理的催化剂在反应条件下吸附和反应，
然后从室温程序升温至所要求的温度，使吸附在催化剂上的各表面物种边反应边脱附出来。
第二种：用作脱附的载气本身就是反应物，在升温过程中载气（或载气中某组分）与催化剂表面
反应形成的吸附物种一面反应一面脱附。两种方式都离不开吸附物种的反应和产物的脱附；
实际上 TPD 和 TPSR 没有严格的区分——判据是升温过程中\emph{是否发生了化学转化}，
而这一点只能靠产物分布（多通道质谱）来确认，不能靠谱图形状来确认。
\end{principle}

\begin{dataflow}
\textbf{实例：CO 在 Ni/\ce{Al2O3} 上的甲烷化。} 吸附均在室温下进行，两种实验方案：
① 将饱和吸附 CO 的催化剂脉冲引入 \ce{H2} 直到不再吸附为止，然后在 He 载气流中做 TPSR；
② 催化剂在含 25\% \ce{H2} 的 \ce{H2}/He 气流中用 CO 饱和吸附，然后在 \ce{H2}--He 气流中做 TPSR。

观测①（共吸附、He 载气）：\ce{CH4} 在 \qty{340}{\degC} 出现一个宽峰；
水峰的起始温度与 \ce{CH4} 基本相同，并有拖尾；此外还能看到 CO 和 \ce{CO2} 的脱附峰。

观测②（流动 \ce{H2}）：\ce{CH4} 在 \qty{225}{\degC} 出现峰；有两个 \ce{H2O} 峰，
低温峰的峰形和峰温与 \ce{CH4} 相同，高温峰则从 \qty{350}{\degC} 开始；
只在 \qty{150}{\degC} 附近有一个宽大的 CO 峰；未见 \ce{CO2} 脱附。

推断：（i）催化剂即使饱和吸附了 CO，也还能吸附大量氢，说明 CO 和 \ce{H2} 吸附在两种不同中心上；
（ii）甲烷在开始生成 \ce{CO2} 的温度下出现；
（iii）在流动氢中，相同温度下以\emph{相同峰形}生成了等量甲烷和水（指第一个水峰），
这意味着甲烷和水的生成速度受\emph{相同}的基元步骤控制。
若表面碳是甲烷化的中间物，则「同时生成甲烷和水」这一速度控制步骤就不可能是这种碳（或氧）
与吸附氢的反应，因为氧与吸附氢的反应速率很快；故速度控制步骤必然是 CO 的 C--O 键断裂。
\end{dataflow}

原文给出的反应步骤为：
（1）\ce{CO(g) <=> CO(a)}；
（2）\ce{CO(a) <=> C(a) + O(a)}；
（3）\ce{CO(a) + O(a) -> CO2(a) -> CO2(g)}；
（4）\ce{2H(a) + O(a) -> H2O(a) -> H2O(g)}；
（5）\ce{C(a) + 4H(a) -> CH4(a) -> CH4(g)}。
\emph{核对说明}：已对照原页面图像（书内 p.339 / PDF p.358）。五步基元反应的箭头类型与化学计量均与原书一致：第（1）（2）步原书用可逆箭头，第（3）（4）（5）步用不可逆箭头。其中第（3）步原书印作 \ce{CO(a) + O(a) -> CO2(a) -> CO2(g)}，正文此前漏去后一段，已补正；其余四步与原书逐字一致，并非重建。

两个方案的差别由此得到解释：对实验①，由于没有足够的氢，反应（3）与反应（5）的速度相当，
所以共吸附实验中能见到 CO、\ce{CO2}、\ce{CH4} 和 \ce{H2O} 的脱附；
对实验②，反应（5）比反应（3）快得多，所以流动氢中只能见到 CO、\ce{CH4} 和 \ce{H2O}。

\begin{pitfall}
\textbf{TPSR 的判据是峰形的对应关系，不是峰温的绝对值。}
上例中最有力的一条证据是「\ce{CH4} 与低温 \ce{H2O} 峰同温度、同峰形」，
它把两个产物绑定到同一个速控步骤上。若只看 \ce{CH4} 峰温（\qty{340}{\degC} 与 \qty{225}{\degC}），
会误以为两种实验方案活化的是不同的活性中心；实际上温度差异来自氢分压不同造成的稳态覆盖度差异，
不是中心种类差异。同理，TPSR 中（5）式的判断依赖于「\ce{CO2} 是否出现」这一\emph{定性}特征，
而 \ce{CO2} 峰的有无很容易被冷阱、色谱柱或质谱扫描速度抹掉，必须事先确认分析系统的响应。
\end{pitfall}

\subsection{误判清单与边界条件}
\label{sec:ch05b-pitfall}

\begin{pitfall}
本节把前面散落的误判集中成一张清单。每一条都给出「错在哪」与「怎么避免」。

\textbf{（1）把 TPD 峰温直接换算成脱附活化能，而不做再吸附与升温速率校正。}
错在把 \equref{eq:ch05b-redhead} 当成恒等式。该式要求一级脱附、$\nu$ 取定、无再吸附、无扩散限制。
多孔体系里再吸附使峰温上移、峰形拖尾，按 Redhead 反算得到的是表观活化能，可偏高数十 kJ\,mol$^{-1}$，
且偏差量随孔结构而变。避免办法：报告时同时给出 $\nu$、$\beta$ 与变速率/变流速实验的结果；
若 $T_{\mathrm p}$ 对流速不收敛，就只报相对强弱，不报 $E_{\mathrm d}$。

\textbf{（2）TPR 耗氢量不扣基线。} 错在把仪器漂移当化学信号。TCD 基线随温度、载气组成和水的陷获而漂移，
不扣基线会把漂移一并积分。避免办法：同条件空白（纯载体）走一遍，或用峰前后基线线性外推，
并写明积分窗口；把基线扣除前后的耗氢量都记下来，二者之差就是这一项的不确定度。

\textbf{（3）把多个重叠峰当成多个物种。} 错在把「可分立的极大值」等同于「独立的化学物种」。
CuO--CeO\textsubscript{2} 中 $\gamma$ 只是 $\beta$ 的肩峰；Ce\textsubscript{0.1}Ti\textsubscript{0.9}O\textsubscript{2}
两峰互相重叠成一个峰。避免办法：分峰拟合 + 变升温速率下各分量是否同步移动 + 独立变量实验
（酸洗、负载量序列、XRD 物相）三选二。

\textbf{（4）氧化态未预处理的 TPR 无意义。} 错在起始价态不确定。避免办法：规定温度的 \ce{O2}/空气预处理，
惰性气氛降温，再切 \ce{H2}/Ar；并把预处理程序写进图题。

\textbf{（5）用热导峰面积代替物种鉴别。} 错在把「脱附物总量」当成「原探针分子量」。
原文明确提示热导信号不能说明脱出物种的化学组成，乙烯/丁烯在分子筛上会聚合，
脱出物的相对分子质量大于原料。避免办法：多通道质谱或冷阱收集 + FID 分析产物分布。

\textbf{（6）跨样品直接比较峰温而不同时比较耗氢量。} 错在把动力学量与容量量混为一谈。
峰又高又陡的样品峰顶温度天然偏低；应按初始耗氢速率法在同一速率下比较温度
（\S\ref{sec:ch05b-quant}），或同时报告耗氢量。

\textbf{（7）把 TPR/TPO 当作非破坏性表征。} 错在忽略样品在程序中已被改变。
氧气 TPD 会把晶格氧抽走并导致结构塌陷与烧结，还原性 TPR 会把样品永久还原。
避免办法：把 TPD/TPR/TPO 排在表征序列的最后，或为每一步预留独立样品。
\end{pitfall}

\begin{crosscheck}
\textbf{三条独立证据链，缺一不可。}（i）谱学量：TPD/TPR/TPO 的峰位、峰面积与峰形；
（ii）结构量：XRD 物相、比表面积、晶粒尺寸；（iii）反应量：转化率、选择性、以活性位计的单位反应速率。
本节各实例中，凡是归属可靠的结论，都是三条链同时成立的——
CuO--CeO\textsubscript{2} 的 $\alpha$/$\beta$/$\gamma$ 归属靠「酸洗 + XRD + 单位反应速率」；
V\textsubscript{2}O\textsubscript{5}/TiO\textsubscript{2} 的物种归属靠「TPR + TPO + XRD」；
Co/\ce{Al2O3} 的积炭中心数靠「\ce{CO2}/\ce{H2O}/\ce{O2} 三通道 + \ce{CH4}/\ce{CO2}-TPD 对照」。
反过来，只靠单一谱图给出的归属（例如仅凭峰温判断 PdO 物种）在原文中也被作者自己标注为「比较困难」。
\end{crosscheck}

\begin{methodsum}
程序升温技术本质上是一种温度谱，其研究对象是特定的探针或反应物分子与催化剂表面特定部位的相互作用，
因此得到的信息对催化研究尤其重要。由于它能获得吸附态分子之间的相互作用及其化学变化，
成为表面反应机理研究的重要手段。许多催化剂特别是氧化催化剂的催化性能往往取决于活性组分的
氧化--还原性质，而程序升温还原和程序升温氧化技术是目前研究催化剂氧化--还原性能最为有效和方便的手段。
使用时须记住三条边界：峰温是动力学量（受升温速率、流速、传热与再吸附影响），
峰面积是容量量（受标定、基线、积分窗口与水陷获影响），峰数是结构量（受峰重叠与吸附质转化影响）；
三者相互独立，任何一条单独使用都会产生系统性的误判。
\end{methodsum}

\srcnote{本分片覆盖上册 \S5.3.6、\S5.4、\S5.5、\S5.6 与 \S5.7，
对应书内 p.320--340 / PDF p.339--359。分片 a 已给出化学吸附原理、三种吸附等温式与 TPD 理论基础，
本分片不重复。文中所有谱图判读规则均取自正文文字描述与变量定义；
凡 OCR 层的失真之处，均已逐条对照原页面图像核验，并在正文以 \emph{核对说明} 记明结论。}
